% GENERATED FILE. Edit canonical lessons, then run npm run generate:books.
% canonical-source-hash: fnv1a64:4f16f734e753241b
% canonical-lessons: LA-C137-faber, LA-C137-pistor, LA-C137-sutor, LA-C137-tonsor, LA-C137-pastor

\chapter{Craftsman, Baker, Shoemaker, Barber, Shepherd}
\label{ch:la-craftsman-baker-shoemaker-barber-shepherd}
\hlchaptermodality{137}

\begin{chapteropening}
\textbf{By the end of this chapter:} \emph{I can say a craftsman, a baker, a shoemaker, a barber and a shepherd.}

Complete the last lesson of chapter 137: i can say a craftsman, a baker, a shoemaker, a barber and a shepherd.
\end{chapteropening}

\section[faber, fabrī]{faber, fabrī — a craftsman, a smith}

\label{lesson:LA-C137-faber}



\begin{quote}
Before the new one: say the Latin for a bone, then the Latin for the stomach.
\end{quote}

\subsection*{You'll want to know: faber, fabrī}
\textbf{faber, fabrī} — \textquotedblleft{}a craftsman, a smith\textquotedblright{}.

\begin{grammarlens}[title={What you've built}]
One more word for this chapter.
\end{grammarlens}

\subsection*{Your turn}
\begin{itemize}
\raggedright
  \item \emph{Say it:} \emph{faber}
  \item \emph{Say it:} \emph{faber}, once more
  \item \emph{Say it:} say \emph{faber}
  \item \emph{Recall:} say the Latin for a bone, then the Latin for the stomach, then say \emph{faber} again
\end{itemize}

\begin{tcolorbox}[breakable,colback=teal!4,colframe=teal!35!black,title={Before you move on}]
What does faber, fabrī mean? (\textquotedblleft{}A craftsman, a smith\textquotedblright{}.) Say it once more.
\end{tcolorbox}
\section[pistor, pistōris]{pistor, pistōris — a baker}

\label{lesson:LA-C137-pistor}



\begin{quote}
Before the new one: say the Latin for the stomach, then the Latin for a craftsman.
\end{quote}

\subsection*{You'll want to know: pistor, pistōris}
\textbf{pistor, pistōris} — \textquotedblleft{}a baker\textquotedblright{}.

\begin{grammarlens}[title={What you've built}]
One more word for this chapter.
\end{grammarlens}

\subsection*{Your turn}
\begin{itemize}
\raggedright
  \item \emph{Say it:} \emph{pistor}
  \item \emph{Say it:} \emph{pistor}, once more
  \item \emph{Say it:} say \emph{pistor}
  \item \emph{Recall:} say the Latin for the stomach, then the Latin for a craftsman, then say \emph{pistor} again
\end{itemize}

\begin{tcolorbox}[breakable,colback=teal!4,colframe=teal!35!black,title={Before you move on}]
What does pistor, pistōris mean? (\textquotedblleft{}A baker\textquotedblright{}.) Say it once more.
\end{tcolorbox}
\section[sūtor, sūtōris]{sūtor, sūtōris — a shoemaker}

\label{lesson:LA-C137-sutor}



\begin{quote}
Before the new one: say the Latin for a craftsman, then the Latin for a baker.
\end{quote}

\subsection*{You'll want to know: sūtor, sūtōris}
\textbf{sūtor, sūtōris} — \textquotedblleft{}a shoemaker\textquotedblright{}.

\begin{grammarlens}[title={What you've built}]
One more word for this chapter.
\end{grammarlens}

\subsection*{Your turn}
\begin{itemize}
\raggedright
  \item \emph{Say it:} \emph{sūtor}
  \item \emph{Say it:} \emph{sūtor}, once more
  \item \emph{Say it:} say \emph{sūtor}
  \item \emph{Recall:} say the Latin for a craftsman, then the Latin for a baker, then say \emph{sūtor} again
\end{itemize}

\begin{tcolorbox}[breakable,colback=teal!4,colframe=teal!35!black,title={Before you move on}]
What does sūtor, sūtōris mean? (\textquotedblleft{}A shoemaker\textquotedblright{}.) Say it once more.
\end{tcolorbox}
\section[tōnsor, tōnsōris]{tōnsor, tōnsōris — a barber}

\label{lesson:LA-C137-tonsor}



\begin{quote}
Before the new one: say the Latin for a baker, then the Latin for a shoemaker.
\end{quote}

\subsection*{You'll want to know: tōnsor, tōnsōris}
\textbf{tōnsor, tōnsōris} — \textquotedblleft{}a barber\textquotedblright{}.

\begin{grammarlens}[title={What you've built}]
One more word for this chapter.
\end{grammarlens}

\subsection*{Your turn}
\begin{itemize}
\raggedright
  \item \emph{Say it:} \emph{tōnsor}
  \item \emph{Say it:} \emph{tōnsor}, once more
  \item \emph{Say it:} say \emph{tōnsor}
  \item \emph{Recall:} say the Latin for a baker, then the Latin for a shoemaker, then say \emph{tōnsor} again
\end{itemize}

\begin{tcolorbox}[breakable,colback=teal!4,colframe=teal!35!black,title={Before you move on}]
What does tōnsor, tōnsōris mean? (\textquotedblleft{}A barber\textquotedblright{}.) Say it once more.
\end{tcolorbox}
\section[pāstor, pāstōris]{pāstor, pāstōris — a shepherd}

\label{lesson:LA-C137-pastor}



\begin{quote}
Before the new one: say the Latin for a shoemaker, then the Latin for a barber.
\end{quote}

\subsection*{You'll want to know: pāstor, pāstōris}
\textbf{pāstor, pāstōris} — \textquotedblleft{}a shepherd\textquotedblright{}.

\begin{grammarlens}[title={What you've built}]
One more word for this chapter.
\end{grammarlens}

\subsection*{Your turn}
\begin{itemize}
\raggedright
  \item \emph{Say it:} \emph{pāstor}
  \item \emph{Say it:} \emph{pāstor}, once more
  \item \emph{Say it:} say \emph{pāstor}
  \item \emph{Recall:} say the Latin for a shoemaker, then the Latin for a barber, then say \emph{pāstor} again
\end{itemize}

\begin{tcolorbox}[breakable,colback=teal!4,colframe=teal!35!black,title={Before you move on}]
What does pāstor, pāstōris mean? (\textquotedblleft{}A shepherd\textquotedblright{}.) Say it once more.
\end{tcolorbox}
